\makeatletter
\def\input@path{{../styles/}{../../styles/}{../../../styles/}{../}{../../}{../../../}}
\makeatother

\documentclass{ee102_pset}
% macros.tex - Course meta information
\renewcommand{\course}{EE 102} % with a space
\renewcommand{\coursetitle}{Signal Processing and Linear Systems}
\renewcommand{\instructor}{Ayush Pandey}
\renewcommand{\student}{Name: }

\renewcommand{\semester}{Fall 2025}
\date{\semester} % this sets the LaTeX date field safely

% Problem set number
\renewcommand{\psetnum}{1}

% Release / due — use renewcommand because package provides empties
\renewcommand{\releasedate}{August 20, 2025}
\renewcommand{\duedate}{August 27, 2025}

\input{preamble.tex}

\author{\rule{3cm}{0.4pt}}
\submitdate{\rule{3cm}{0.4pt}}
\problemset{Problem Set \#12: Sampling}
\renewcommand{\psetnum}{12}
\renewcommand{\releasedate}{November 18, 2026}
\renewcommand{\duedate}{Dec 2, 2026 at 11:59 PM}
\shorttitle{Problem Set \#12}
\hypersetup{pdftitle={EE 102 Problem Set 12: Sampling}}

\begin{document}

% Fall 2025 homework/week12/hw11.tex, Problem 2(a)--(e).
\problem{1} [Adapted from Problem 7.2 in Oppenheim \& Willsky]

A signal $x(t)$ with Fourier transform $X(j\omega)$ undergoes impulse-train sampling (like what we discussed during the lectures) to generate
\[
x_p(t) = \sum_{n=-\infty}^{\infty} x(nT)\,\delta(t-nT)
\]
where $T = 10^{-4}$. For each of the following sets of constraints on $x(t)$ and/or $X(j\omega)$, does the sampling theorem guarantee
that $x(t)$ can be recovered exactly from $x_p(t)$? Clearly justify your answers.

\textbf{[12 points each]}

\problempart $X(j\omega) = 0$ for $\lvert \omega \rvert > 5000\pi$
\problempart $X(j\omega) = 0$ for $\lvert \omega \rvert > 15000\pi$
\problempart $\Re\{X(j\omega)\} = 0$ for $\lvert \omega \rvert > 5000\pi$
\problempart $x(t)$ real and $X(j\omega) = 0$ for $\omega < -5000\pi$
\problempart $x(t)$ real and $X(j\omega) = 0$ for $\omega < -15000\pi$

% Fall 2025 homework/week12/hw11.tex, Problem 3(a).
\problem{2} [Adapted from Problem 7.22 in Oppenheim \& Willsky]

The signal $y(t)$ is generated by convolving a band-limited signal
$x_1(t)$ with another band-limited signal $x_2(t)$, that is,
\[
y(t) = x_1(t) * x_2(t)
\]
where
\[
X_1(j\omega) = 0 \quad \text{for } \lvert \omega \rvert > 1000\pi,
\qquad
X_2(j\omega) = 0 \quad \text{for } \lvert \omega \rvert > 2000\pi.
\]

Impulse-train sampling is performed on $y(t)$ to obtain
\[
y_p(t) = \sum_{n=-\infty}^{\infty} y(nT)\,\delta(t-nT).
\]

\problempart \textbf{[40 points]} Specify the range of values for the sampling period $T$ which ensures that
$y(t)$ is recoverable from $y_p(t)$.

\end{document}
